% Replaces the subsection "The D4 Group Structure of Hohfeldian Transitions" in section 8.4
% (Geometric-Ethics v1.24) at the LaTeX conversion. Ledger: GE-PROP-8-3-V4 (proved), which
% supersedes GE-PROP-8-3; GE-D4-MEASURED withdrawn. Owner decision 2026-10-03: D4 is obsolete.

\subsubsection{The $V_4$ Structure of Hohfeldian Transitions}\label{sec:hohfeld-v4}

The four Hohfeldian jural positions, Obligation (O), Claim (C), Liberty (L) and No-claim (N), are
related by two independent binary oppositions:
\begin{itemize}
  \item the \emph{correlative} opposition $s$: $O \leftrightarrow C$, $L \leftrightarrow N$, the
        same relation seen from the other party's side;
  \item the \emph{jural} opposition $n$: $O \leftrightarrow L$, $C \leftrightarrow N$, deontic
        negation.
\end{itemize}

\begin{proposition}[$V_4$ structure of Hohfeldian transitions]\label{GE-PROP-8-3-V4}
The operations $s$ and $n$ are commuting involutions, and the group they generate acting on the
four positions is the Klein four-group
\[
  V_4 = \{e,\; n,\; s,\; sn\} \;\cong\; \mathbb{Z}_2 \times \mathbb{Z}_2 ,
\]
abelian, with every element its own inverse. The action is regular: for any two positions there is
exactly one operation taking the first to the second. In particular no operation cycles the four
positions, so the quarter-turn $O \to C \to N \to L \to O$ lies outside the group.
\end{proposition}

\noindent\emph{Status: proved.} Machine-checked in Lean~4 with Mathlib (erisml-lib,
\texttt{formal/HohfeldV4.lean}: commuting involutions, \texttt{card\_closure\_eq\_four},
\texttt{measured\_sector\_abelian}, \texttt{quarter\_turn\_not\_generated}).

The encoding that makes the structure transparent writes each position as two bits, whether it is
negated and whether it is correlated, relative to Obligation, and each operation as the bits it
flips:
\[
  O = 00,\quad L = 10,\quad C = 01,\quad N = 11; \qquad e = 00,\quad n = 10,\quad s = 01,\quad sn = 11 .
\]
Applying an operation to a position, and composing two operations, are then both bitwise XOR.

\paragraph{What was claimed before, and why it is withdrawn.} Earlier versions of this book stated
that the transitions form the dihedral group $D_4$ of order eight, the symmetry group of the
square, adding a quarter-turn $r$ that cycles the four positions. Three facts retire that claim.
No quarter-turn has ever been demonstrated as a normative operation. The Lean proof shows the
quarter-turn is not generated by the two operations that are demonstrated. And a direct search
for it in a learned representation of normative statements did not find it (erisml-lib,
\texttt{docs/papers/quarter-turn-hunt}): the fitted cycle map lacked the non-abelian signature.
The measured correlative symmetry rates ($O \leftrightarrow C$ at 87\%, $L \leftrightarrow N$ at
82\%) support the two involutions, which is to say $V_4$, and nothing more. The corollary that
``a double promise creates a claim'' rested on $r^3 \cdot r^3 = r^2$ and is withdrawn with it.

\paragraph{Semantic gates as $V_4$ elements.} A semantic gate (Definition~8.8) acts from its
source stratum, so each gate is one $V_4$ operation restricted to the positions it applies from:
\begin{itemize}
  \item ``You promised'' takes $L \to O$, and ``only if convenient'' takes $O \to L$: both are the
        jural flip $n$.
  \item ``From their perspective'' is the correlative $s$.
  \item ``You have every right'' takes $L \to C$, and ``they can't demand'' takes $C \to L$: both are
        $sn$, the unique operation between those positions.
\end{itemize}
Because a gate fires only from its source stratum, a second promise made in the obligation stratum
changes nothing, as it should.

\paragraph{Generator asymmetry.} [Empirical result (preliminary).] The asymmetry between the two
generators survives the correction intact; only the name of the gated generator changes. The
correlative $s$ is non-gated, immediate and free: any competent reasoner can adopt the other
party's perspective at will, which explains the high correlative symmetry observed in
\S17.3. The jural flip $n$ is discrete and gated: it needs a trigger (``you promised'', ``only if
convenient'', ``he hit you'') and fires with probability 89--97\% (the gate detection table above).
The first is a change of frame; the second is a change of state, which requires an event.

\paragraph{Order independence, and what an order effect means.} $V_4$ is abelian, so the holonomy of
any path of operations depends only on how many times $s$ and $n$ occur (mod~2), never on their
order. A reasoner whose verdicts depend on the order of a perspective shift and a negation is
therefore violating the structure, and the violation is measurable: the bond index and the Wilson
observable of the reference implementation (erisml-lib \texttt{erisml.ethics.hohfeld}, V4 since
October 2026) measure exactly this.
